\section{Introduction}

\subsection{Project overview}
Rùa Toán Basic is a mathematics practice platform designed for lower secondary students. It represents
a deliberate, ground-up rebuild of the Rùa Toán learning platform with a focus on simplicity,
architectural transparency, and pedagogical alignment. The frontend is implemented entirely with
plain HTML5, CSS3, and vanilla JavaScript without any framework dependencies or build steps; the backend
is powered by a lightweight FastAPI service running serverless on Vercel; user accounts and homework photo
submissions are managed cleanly in PostgreSQL. Both versions share the exact same comprehensive curriculum
catalogue (3,381 collections, 33,929 problems, 90 lessons, 216 books, and 89 documents) stored as versioned
JSON files.

The live production deployment for this basic version is hosted at
\texttt{https://ruatoanbasic.vercel.app} (source code: \texttt{https://github.com/thunopro/ruatoan\_basic}).
The original flagship version is available at \texttt{https://ruatoan.vercel.app}
(source code: \texttt{https://github.com/thunopro/ruatoan}).

\subsection{Motivation and pedagogical context: from the advanced platform to the basic foundation}
The creation of Rùa Toán Basic arose from an honest reflection on the learning journey in web development.
Initially, the author conceived and ``vibed'' the flagship Rùa Toán platform using an advanced, state-of-the-art
technology stack: Next.js 15, React 19, TypeScript, Tailwind CSS, and Auth.js. While that version was ambitious
and feature-rich, it relied heavily on framework abstractions---server-side rendering pipelines, complex
component hydration, intricate bundler magic, and automated authentication middleware.

For a student beginning with the university curriculum of \textit{Web Design and Programming}, that high level of
abstraction proved overwhelming to fully grasp from first principles. It masked the fundamental mechanics of the web:
how an HTTP request is structured, how the browser DOM is updated imperatively, how cookies are cryptographically signed,
and how client scripts communicate with a RESTful backend API.

Recognizing that true mastery requires understanding the foundations, the author made the deliberate decision to
pivot and ``vibe'' a fresh, transparent version: \textbf{Rùa Toán Basic}. This rebuild intentionally eliminates
all framework scaffolding in order to create a genuine environment for learning, experimentation, and mastery:
\begin{enumerate}
  \item \textbf{Honest learning curve and inspectability.} By rebuilding the interface in vanilla HTML/CSS/JS and
        the backend in straightforward FastAPI with direct SQL queries, every layer of the web application becomes
        transparent. A student can open browser DevTools, inspect the actual network payloads and DOM nodes,
        and follow the data flow directly to PostgreSQL without cognitive framework overhead.
  \item \textbf{Pedagogical curriculum alignment.} The basic version directly aligns with the learning outcomes
        of the course: mastering semantic HTML, responsive CSS layouts (flexbox, grid, custom properties),
        JavaScript event handling, RESTful API design, and cookie-based session security.
  \item \textbf{Architectural clarity and controlled scope.} Stripping away opaque build pipelines allowed the author
        to focus deeply on core platform capabilities: rendering LaTeX mathematics with MathJax, implementing
        PBKDF2-HMAC password hashing, validating image magic bytes for homework uploads, and deploying reliably
        to serverless edge infrastructure.
\end{enumerate}

\subsection{Objectives}
The basic version was built to satisfy the following objectives:

\begin{enumerate}[label=O\arabic*.]
  \item Provide the core learning loop: browse exam collections, read individual problems
        with mathematical notation, study lessons, and consult books and documents.
  \item Support filtering and search: by grade, topic, collection kind, year, province, and
        free text, with pagination.
  \item Render mathematics correctly in the browser using MathJax, including in problem
        statements and lessons.
  \item Allow a learner to register, log in, and log out with an email and password, with
        the session held in a signed cookie.
  \item Allow a logged-in learner to upload, view, and delete photographs of written work,
        with the images stored in PostgreSQL.
  \item Let a learner mark problems in the browser and keep that list locally, without
        requiring an account.
  \item Deploy successfully on Vercel, with durable storage in PostgreSQL rather than in
        the function's temporary filesystem.
  \item Keep the source small, readable, and free of a build step.
\end{enumerate}

\subsection{Scope and boundaries}
Implemented in this version:

\begin{itemize}
  \item Home page with live catalogue totals and a list of marked problems.
  \item Practice page: text/grade/topic filters, 20 results per page, problem cards, and a
        ``mark / unmark'' toggle.
  \item Exam collections page: paginated list (20 per page) and a collection detail page
        showing problems in their original order, with source attribution and a link to
        the source.
  \item Lessons page: client-side search over 90 lessons, expandable bodies.
  \item Library page: books (216) and documents (89) with client-side search and external
        links.
  \item Account pages: registration and login with email and password; logout.
  \item Submission page: upload up to six images (1.5 MB each) per request, view them,
        and delete them.
  \item Catalogue of 3,381 collections and 33,929 problems, unchanged from the original.
\end{itemize}

Explicitly \emph{not} implemented in this version, by design:

\begin{itemize}
  \item Google OAuth / social sign-in (the original has it; this version uses email and
        password only).
  \item Device-token based pseudonymous photo ownership (this version ties photos to the
        logged-in account).
  \item Cross-device synchronisation of marked problems (they stay in \texttt{localStorage}).
  \item Automatic answer marking, a teacher administration panel, or a learning analytics
        dashboard.
  \item The content crawl and preparation pipeline (the JSON is reused as-is).
\end{itemize}

\subsection{Technology overview}
The stack is deliberately small. Table~\ref{tab:technology} lists each layer, the
technology, and its responsibility in Rùa Toán Basic.

\begin{table}[ht]
\centering\small
\begin{tabularx}{\textwidth}{@{}l l X@{}}
\toprule
\textbf{Layer} & \textbf{Technology} & \textbf{Use in Rùa Toán Basic}\\
\midrule
Markup & Plain HTML5 & Eight pages (home, practice, exams, lesson, library, login,
        register, submissions); semantic markup, no template engine\\
Styling & Plain CSS3 & Unified styling: \texttt{public/css/original.css} (core visual layout,
        branding, cards, footer) + \texttt{public/css/style.css} (design tokens, forms,
        photo grid, a11y, dark mode)\\
Behaviour & Plain JavaScript (ES2020) & Eight small files (\texttt{public/js/*.js}); modular,
        no framework, no bundler, no build step\\
Mathematics & MathJax 3 (CDN) & Client-side LaTeX rendering of inline \texttt{\$...\$}
        and display \texttt{\$\$...\$\$}\\
Backend & FastAPI + Uvicorn & \texttt{backend/main.py} and helpers \texttt{auth.py},
        \texttt{content.py}, \texttt{database.py}; about 600 clean lines\\
Database & PostgreSQL via \texttt{psycopg} & Schema \texttt{ruatoan\_basic}: tables
        \texttt{users}, \texttt{photos}, and \texttt{password\_resets}\\
Catalogue & JSON files & 3,381 collections, 33,929 problems, 90 lessons, 216 books/resources,
        89 documents\\
Hosting & Vercel (Python runtime) & Root \texttt{main.py} entry point, Vercel function
        routing, automated CI/CD\\
\bottomrule
\end{tabularx}
\caption{Technologies used in Rùa Toán Basic.}
\label{tab:technology}
\end{table}

\subsection{Catalogue size}
The content catalogue is reused unchanged from the original project. The inspected JSON
files contain 3,381 collection records and 33,929 nested problem records, together with
90 lessons, 216 book, magazine, or resource records, and 89 document records. That is
37,705 raw records when each collection, problem, lesson, book, and document entry is
counted once. The JSON files occupy roughly 16 MB on disk (after removing a 19 MB
duplicate copy, see Section~\ref{sec:cleanup}). A problem may carry more than one topic
tag, so topic counts are assignments and do not partition the problem set.

\subsection{Report evidence and snapshot}
This report is based on a source-code audit and live verification of the Rùa Toán Basic deployment
on 1 November 2026, following architectural unification, security hardening, and accessibility
fixes. The following checks were run and verified:

\begin{itemize}
  \item \texttt{node --check} on all eight JavaScript files (all passed).
  \item \texttt{python3 -m py\_compile} on all Python backend and script files (all passed).
  \item Automated smoke test across all routes (\texttt{/}, \texttt{/api/stats},
        \texttt{/api/daily}, \texttt{/api/collections},\\
        \texttt{/api/problems}, \texttt{/api/lessons},
        \texttt{/api/books}, and \texttt{/api/docs}): each returned HTTP 200 with valid payloads.
  \item Complete account and security flow: user registration, credential validation, session cookie
        issuance, authenticated \texttt{/api/me}, password change, database-persisted password reset
        via single-use tokens, and logout all executed end to end.
  \item Photo submission pipeline: magic-byte validation (PNG/JPEG/WebP), size bounds (1.5 MB),
        database binary storage, retrieval with \texttt{nosniff}, and ownership-verified deletion.
  \item Browser runtime audit: zero console errors across all eight pages in production, zero
        404 asset failures (fonts and favicons fully resolved), and validated keyboard accessibility
        (\texttt{:focus-visible}) and high-contrast inputs.
\end{itemize}

Database credentials and session secrets are managed strictly through environment variables
(\texttt{.env.local} and Vercel environment settings) and never committed to source control.
